\chapter{Weiss and Monneau monotonicity}\label{sec:obstacle-monotonicity}
Monotonicity identifies the homogeneity of blow-ups and removes subsequence ambiguity at singular points. Throughout this section $u$ is a normalized obstacle solution as in Definition \ref{def:normalized-obstacle}, the chosen free-boundary point has been translated to zero, and $\overline{B_R}\subset U$. We follow \cite{Fig18}, Section 8, pp.~17--19. Derivative identities are stated almost everywhere in the radius, with local absolute continuity supplying monotonicity at every radius.

\begin{definition}[The Weiss energy]\label{def:obstacle-weiss-energy}
For $0<r<R$, define
\[
 W(r,u)=r^{-n-2}\int_{B_r}(|Du|^2+2u)\,dx
       -2r^{-n-3}\int_{\partial B_r}u^2\,d\mathcal H^{n-1}.
\]
The coefficient $2$ in the bulk term corresponds to the equation $\Delta u=1$ on positivity, and the boundary coefficient corresponds to quadratic scaling. Direct change of variables gives $W(r,u)=W(1,u_{0,r})$. This is \cite{Fig18}, Proposition 8.2.
\end{definition}

\begin{theorem}[Weiss monotonicity and its exact derivative]\label{thm:obstacle-weiss-monotonicity}
Let $u$ be a normalized obstacle solution with $0\in\Gamma$ and $\overline{B_R}\subset U$. The function $r\mapsto W(r,u)$ is locally absolutely continuous and
\[
 \frac{d}{dr}W(r,u)
 =\frac{2}{r^{n+2}}\int_{\partial B_r}
 \left(\partial_\nu u-\frac{2u}{r}\right)^2\,d\mathcal H^{n-1}
 \geq0
\]
for almost every $r$. It has a finite limit $W(0+,u)$ as $r\downarrow0$.
\end{theorem}
\begin{proof}
Write $u_r(x)=r^{-2}u(rx)$. On a compact interval of positive radii, the $C^{1,1}$ bounds justify differentiation under the fixed-ball integrals and integration by parts for almost every radius. In particular $W(1,u_r)$ is absolutely continuous and
\[
 \frac{d}{dr}W(1,u_r)
 =2\int_{B_1}(1-\Delta u_r)\partial_ru_r
  +2\int_{\partial B_1}(\partial_\nu u_r-2u_r)\partial_ru_r.
\]
The interior integrand vanishes: where $u_r>0$ the first factor is zero, while on contact both $u_r$ and $Du_r$ vanish and
$\partial_ru_r=r^{-1}(x\cdot Du_r-2u_r)=0$. On $\partial B_1$ this last identity becomes $r^{-1}(\partial_\nu u_r-2u_r)$. Hence
\[
 \frac{d}{dr}W(1,u_r)=\frac2r\int_{\partial B_1}(\partial_\nu u_r-2u_r)^2.
\]
Scaling the surface integral gives the displayed formula. Quadratic upper growth and $|Du(x)|\leq L|x|$ bound both terms in $W$ uniformly at small radii, so its monotone limit is finite. This is the computation in \cite{Fig18}, Proposition 8.2, with its regularity justification and scaling factors explicit.
\end{proof}

\begin{corollary}[Equality in Weiss monotonicity]\label{cor:obstacle-weiss-equality}
For $u$ as in Theorem \ref{thm:obstacle-weiss-monotonicity}, the Weiss energy is constant on an interval $(a,b)\subset(0,R)$ if and only if
\[
 x\cdot Du(x)=2u(x)\quad(a<|x|<b).
\]
Equivalently, along every ray the quotient $u(r\omega)/r^2$ is constant while $a<r<b$. Constancy on $(0,b)$ means that $u$ is homogeneous of degree two in $B_b$.
\end{corollary}
\begin{proof}
If the energy is constant, the nonnegative derivative formula makes $x\cdot Du-2u$ vanish on almost every sphere in the annulus. Coarea and continuity make it vanish everywhere there. Along a ray,
$\frac d{dr}(r^{-2}u(r\omega))=r^{-3}(r\omega\cdot Du-2u)$, which proves the homogeneity assertion. Conversely that identity makes the derivative of $W$ zero almost everywhere; absolute continuity makes $W$ constant. This gives the equality case of Theorem \ref{thm:obstacle-weiss-monotonicity}.
\end{proof}

\begin{proposition}[Every obstacle blow-up is quadratic homogeneous]\label{prop:obstacle-blowup-homogeneous}
Every blow-up $u_0$ in Theorem \ref{thm:obstacle-blowup-compactness} satisfies
\[
 u_0(tx)=t^2u_0(x)\quad(t>0),\qquad
 W(s,u_0)=W(0+,u)\quad(s>0).
\]
\end{proposition}
\begin{proof}
If $u_{r_j}\to u_0$ in $C^1_{\mathrm{loc}}$, then for any fixed $s>0$ the bulk and boundary integrals converge, giving
$W(s,u_0)=\lim_jW(s,u_{r_j})=\lim_jW(sr_j,u)=W(0+,u)$. The limit is a normalized obstacle solution, so the equality case on every interval of radii proves global degree-two homogeneity. This uses no classification of homogeneous solutions.
\end{proof}

\begin{proposition}[Weiss energies of the two model families]\label{prop:obstacle-model-weiss}
For the quadratic polynomials $p_A$ and half-parabolas $h_e$,
\[
 W(r,p_A)=\frac{\omega_n}{2(n+2)},\qquad
 W(r,h_e)=\frac{\omega_n}{4(n+2)}
 \quad(r>0).
\]
Thus the limiting Weiss energy distinguishes regular and singular points once the dichotomy theorem is applied.
\end{proposition}
\begin{proof}
Both models are degree-two homogeneous and satisfy $u\Delta u=u$ almost everywhere. Integration by parts on $B_1$ therefore gives
$\int_{B_1}|Du|^2=2\int_{\partial B_1}u^2-\int_{B_1}u$, so $W(1,u)=\int_{B_1}u$. Rotational symmetry gives $\int_{B_1}x_i^2=\omega_n/(n+2)$. Taking the trace of $A$ gives the first value; restriction to a half-ball gives the second. Scaling gives all radii, and $C^1$ blow-up convergence gives the corresponding limiting energies. The common quadratic value is the constant used in \cite{Fig18}, Proposition 8.2.
\end{proof}

\begin{definition}[The Monneau distance to a quadratic model]\label{def:obstacle-monneau}
For $p=p_A$ as in Example \ref{ex:obstacle-quadratic-model}, set
\[
 M(r,u,p)=r^{-n-3}\int_{\partial B_r}(u-p)^2\,d\mathcal H^{n-1}
 =\int_{\partial B_1}(u_r-p)^2\,d\mathcal H^{n-1}.
\]
This is a spherical $L^2$ distance after quadratic rescaling, not an integral over the unit ball. Its monotonicity below requires a singular center. See \cite{Fig18}, Lemma 8.4.
\end{definition}

\begin{theorem}[Monneau monotonicity at a singular point]\label{thm:obstacle-monneau-monotonicity}
Let $u$ be a normalized obstacle solution with zero singular and $\overline{B_R}\subset U$. For every quadratic model $p=p_A$, the function $M(r,u,p)$ is locally absolutely continuous and nondecreasing. More precisely, with $w=u-p$,
\[
 M'(r)=\frac2r\bigl(W(r,u)-W(r,p)\bigr)
       +\frac{2}{r^{n+3}}\int_{B_r}p\chi_{\{u=0\}}\,dx
 \geq0
\]
for almost every $r\in(0,R)$.
\end{theorem}
\begin{proof}
Existence of one quadratic blow-up and Proposition \ref{prop:obstacle-model-weiss} imply
$W(0+,u)=W(r,p)=\omega_n/(2(n+2))$. Thus the first term in the claimed derivative is nonnegative by Weiss monotonicity.

Because $\Delta p=1$ and $x\cdot Dp=2p$, expansion of the energy and integration by parts cancel the cross terms, leaving
\[
 W(r,u)-W(r,p)=r^{-n-2}\int_{B_r}|Dw|^2
                         -2r^{-n-3}\int_{\partial B_r}w^2.
\]
Also $\Delta w=-\chi_{\{u=0\}}$ and $w=-p$ on contact, so $w\Delta w=p\chi_{\{u=0\}}\geq0$. Another integration by parts yields
\[
 r^{-n-3}\int_{\partial B_r}w(x\cdot Dw-2w)
 =W(r,u)-W(r,p)+r^{-n-2}\int_{B_r}w\Delta w.
\]
Differentiating the fixed-sphere formula for $M$ gives
$M'(r)=2r^{-n-4}\int_{\partial B_r}w(x\cdot Dw-2w)$. Combining the last two identities proves the formula and its sign. Local absolute continuity follows from $C^{1,1}$ regularity on compact annuli. This is the argument of \cite{Fig18}, Lemma 8.4, with the nonnegative terms retained.
\end{proof}

\begin{theorem}[Uniqueness and continuity of singular blow-ups]\label{thm:obstacle-singular-uniqueness}
At each singular point $x_0$ of a normalized obstacle solution there is one quadratic model $p_{x_0}$ such that
\[
 u_{x_0,r}\longrightarrow p_{x_0}\quad\text{in }C^1_{\mathrm{loc}}(\mathbb R^n)
 \quad\text{as }r\downarrow0.
\]
Its matrix depends continuously on $x_0$ within the singular set, and uniformly continuously on compact subsets of that set.
\end{theorem}
\begin{proof}
Translate $x_0$ to zero and choose a sequence with quadratic limit $p$. Then $M(r_j,u,p)\to0$. Since this distance is nonnegative and nondecreasing with the radius, its limit at zero is zero. Every other blow-up $q$ consequently agrees with $p$ on $\partial B_1$ in $L^2$, hence everywhere on that sphere by continuity. Both are homogeneous of degree two by Proposition \ref{prop:obstacle-blowup-homogeneous}, so they agree throughout space. Compactness and uniqueness of every subsequential limit give convergence along all radii.

For continuity, take singular points $x_j\to x_0$ and choose a fixed small $\rho>0$ for which all the relevant balls remain in $U$. Apply Monneau monotonicity at $x_j$ with the comparison polynomial $p_{x_0}$. Letting the smaller radius go to zero gives
\[
 \int_{\partial B_1}|p_{x_j}-p_{x_0}|^2
 \leq\int_{\partial B_1}
 \left|\rho^{-2}u(x_j+\rho x)-p_{x_0}(x)\right|^2.
\]
First send $j$ to infinity using continuity of $u$, then send $\rho$ to zero using the just-proved convergence at $x_0$. The resulting spherical $L^2$ distance tends to zero. On the finite-dimensional space of quadratic forms this is a norm, so the matrices converge. Continuity on a compact subset is uniform. This proves the uniqueness and continuity assertions of \cite{Fig18}, Theorem 8.1, using spherical integrals throughout.
\end{proof}

\begin{definition}[Singular strata of the classical obstacle problem]\label{def:obstacle-singular-strata}
Write $p_{x_0}(x)=\frac12x\cdot A_{x_0}x$ for the unique singular blow-up and define
\[
 \Sigma_k=\{x_0\in\Sigma:\dim\ker A_{x_0}=k\},
 \qquad k=0,\ldots,n-1.
\]
The trace-one condition excludes the zero matrix. The kernel is the contact set of the blow-up; its dimension defines the stratum, not the dimension of the full ambient free boundary. See \cite{Fig18}, Section 9, p.~19.
\end{definition}

\begin{theorem}[Classical stratification of singular obstacle points]\label{thm:obstacle-singular-stratification}
For a normalized obstacle solution and each $k=0,\ldots,n-1$, the stratum $\Sigma_k$ is locally contained in a $k$-dimensional $C^1$ submanifold. In particular each point of $\Sigma_0$ is isolated within $\Sigma_0$. Containment does not assert that the whole stratum is itself a manifold.

This is \cite{Fig18}, Theorem 9.1, p.~19, attributed there to Caffarelli. The extension and implicit-function arguments establishing manifold containment remain external. The uniqueness theorem alone does not provide this geometric conclusion.
\end{theorem}
